% Оформление основано на report_pattern курса MAI-IR.
% Уведомление об авторских правах шаблона: styles/template-license.txt.
\documentclass[12pt,a4paper,oneside,fleqn]{article}
\usepackage{styles/log-style}
\usepackage{booktabs,tabularx,caption,float}
\usepackage{pgfplots}
\pgfplotsset{compat=1.18}
\usepackage{microtype}
\hypersetup{unicode=true,hidelinks,pdftitle={Отчёт по информационному поиску},pdfauthor={Книга Тимофей Сергеевич}}
\setlength{\emergencystretch}{3em}
\setcounter{tocdepth}{2}
\renewcommand{\thesubsection}{\thesection.\arabic{subsection}}
\renewcommand{\arraystretch}{1.18}
\lstset{basicstyle=\small\ttfamily,numbers=none,frame=single,breaklines=true,showstringspaces=false,escapechar={},language={},columns=fullflexible}
\newcommand{\code}[1]{\texttt{\detokenize{#1}}}
\newcommand{\lab}[2]{\clearpage\section{Лабораторная работа №#1 «#2»}}
\newcommand{\task}{\subsection*{Задание}}
\newcommand{\method}{\subsection{Описание метода решения}}
\newcommand{\journal}{\subsection{Журнал выполнения}}
\newcommand{\testing}{\subsection{План тестирования и автотесты}}
\newcommand{\results}{\subsection{Результаты}}
\newcommand{\conclusions}{\subsection{Выводы}}
% Данные титульного листа; при необходимости редактируются здесь.
\newcommand{\StudentName}{Книга Тимофей Сергеевич}
\newcommand{\StudentGroup}{М8О-409Б-23}
\newcommand{\TeacherName}{А.\,А. Кухтичев}
\newcommand{\ReportYear}{2026}

\begin{document}
\begin{titlepage}
\begin{center}
\bfseries
{\Large Московский авиационный институт\\
(национальный исследовательский университет)}
\vspace{36pt}

{\large Факультет информационных технологий\\и прикладной математики}
\vspace{30pt}

{\large Кафедра вычислительной математики\\и программирования}
\vspace{42pt}

Отчёт по лабораторным работам\\по курсу «Информационный поиск»
\vspace{18pt}

«Учебная поисковая система по корпусу\\документации и обсуждений программирования»
\vspace{12pt}

Лабораторные работы №1--8, 10--13, 17
\end{center}
\vspace{36pt}
\begin{flushright}
\begin{tabular}{rl}
Студент: & \StudentName \\
Преподаватель: & \TeacherName \\
Группа: & \StudentGroup \\
Дата: & \rule{45mm}{0.3pt} \\
Оценка: & \rule{45mm}{0.3pt} \\
Подпись: & \rule{45mm}{0.3pt}
\end{tabular}
\end{flushright}
\vfill
\begin{center}\bfseries Москва, \ReportYear\end{center}
\end{titlepage}

\tableofcontents
\clearpage
\section*{Общие сведения}
\addcontentsline{toc}{section}{Общие сведения}
В работе последовательно реализованы сбор корпуса, поисковый робот,
токенизация, морфологическая нормализация, оценка качества, построение
обратного индекса, булев и цитатный поиск, сжатие, прыжки по спискам
постингов и ранжирование TF-IDF. Корпус содержит документацию и обсуждения
по программированию, операционным системам, компиляторам и системным
библиотекам. Основные компоненты написаны на C++20, обвязка и робот на Python.

Нумерация соответствует списку лабораторных курса «Информационный поиск»
\cite{course}. Выполнены работы №1--8, 10--13 и 17.
Каждый раздел содержит задание, описание метода, журнал выполнения,
план тестирования, результаты и выводы.

Корпус содержит 1\,679\,971 документ и 17\,236\,388\,223 байта текста.
Для опыта с законом Ципфа отдельно выделены первые 10\,000 документов.
Это техническая выборка, а не случайная репрезентативная выборка всего корпуса.
Размеры, статистика токенизации и индексации относятся к указанному
набору данных; сравнение сжатия и прыжков выполняется при одинаковой
нормализации документов и запросов.

Качество поиска оценено на 30 запросах для четырёх систем.
Размечены все 426 уникальных пар «запрос, документ» объединённого пула
первых пяти результатов. Оценки выставлены ассистентом по содержимому
документов и информационным потребностям запросов, без независимых
человеческих оценщиков. Поэтому метрики характеризуют данный эксперимент
и не являются исчерпывающей оценкой всего корпуса.

\lab{1}{Добыча корпуса документов}
\task
Подготовить корпус для следующих лабораторных: получить документы не менее
чем из двух источников, определить единицу документа, изучить разметку и
метаданные, выделить текст. Найти существующие средства поиска по источникам,
привести примеры запросов и ограничения выдачи. Определить количество
документов, объёмы исходных данных и текста, средние размеры документов.

\method
Корпус объединяет техническую документацию и архивные обсуждения Stack
Exchange. Для документации единицей является HTML-страница с собственным
URL и заголовком; для обсуждений - архивная ветка вопроса. Такое разделение
сохраняет контекст задачи и ответов, но приводит к неодинаковой длине
документов. Источники включают Linux Kernel, man7, cppreference, проекты
LLVM, GNU, FreeBSD и восемь сайтов Stack Exchange.

Документация поступает из экспортированной базы робота. Для большого
числа обсуждений используется набор \code{common-pile/stackexchange}
с Hugging Face \cite{dataset}. Зафиксирована ревизия
\code{5ec3aa2e65b7c61b7c18f1ff18ccbc4d3279a23c};
происхождение и авторство не подменяются результатом собственного обхода.
Изначально порог длины обсуждения составлял 1000 слов; начиная с архива
\code{00024_se.jsonl.gz} для Stack Overflow он снижен до 850 слов.
Страницы документации этим порогом не ограничены.

HTML обрабатывается профилем конкретного сайта: удаляются оформление,
скрипты и навигация, сохраняется основной текст, включая примеры кода.
Для повторного извлечения в исходном экспорте остаются HTML-файлы.
Большой корпус хранится компактно, без отдельного текстового файла на
каждый документ. Формат \code{corpus.bin} начинается сигнатурой
\code{IRDOC001}, далее следуют записи:
\begin{center}
\begin{tabular}{lll}\toprule
Поле & Размер & Содержание\\\midrule
\code{url_length} & 4 байта & Длина URL в байтах\\
\code{title_length} & 4 байта & Длина заголовка в байтах\\
\code{text_length} & 8 байт & Длина текста в байтах\\
URL, заголовок, текст & переменный & UTF-8, подряд без разделителей\\\bottomrule
\end{tabular}
\end{center}
Целые записываются little-endian. Манифест JSONL хранит идентификаторы,
источник, URL, заголовок и смещение текста; сводная статистика хранится
отдельно. Это позволяет читать корпус последовательно и проверять
целостность независимо от индексатора.

\journal
Сначала реализован экспорт из MongoDB и извлечение текста по профилям
сайтов. Затем добавлены импорт архивных обсуждений, проверка архивов и
контрольные точки: прерванный импорт продолжается с последнего сохранённого
состояния. Большое количество мелких файлов заменено упакованным потоком.
Для защиты от неполного экспорта анализатор сверяет смещения и объёмы
с манифестом. Главная проблема неоднородного корпуса - разные понятия
«сырого документа»: для HTML это исходная страница, для архива -
распакованная JSON-строка. В статистике эта договорённость записана явно.

\testing
Проверяются сигнатура потока, длины полей, отсутствие NUL внутри текста,
корректность UTF-8, монотонность смещений и выход за границы файла.
Число записей и сумма длин текстов должны совпасть с итоговыми агрегатами.
Отдельно следует проверять извлечение на страницах разных источников:
осмысленный заголовок, наличие основного текста, сохранение примеров кода,
отсутствие меню. Проверка количества документов сама по себе не выявляет
дублирующиеся тексты, поэтому дедупликация по содержимому остаётся отдельной
задачей. Статистика включает число документов и объёмы
по каждому источнику, что позволяет независимо сверять итоговые суммы.

\results
\begin{table}[H]\centering
\begin{tabular}{lr}\toprule
Показатель & Значение\\\midrule
Количество документов & 1\,679\,971\\
Число источников в статистике & 22\\
Исходные данные, байт & 20\,327\,391\,393\\
Извлечённый текст, байт & 17\,236\,388\,223\\
Упакованный корпус, байт & 17\,433\,701\,772\\
Слов в статистике экспорта & 2\,237\,727\,857\\
Средний исходный размер, байт & 12\,099,85\\
Средний объём текста, байт & 10\,259,93\\
Среднее число слов & 1332,00\\\bottomrule
\end{tabular}
\caption{Статистика основного корпуса. «Слова» экспортёра и токены индексатора считаются разными правилами.}
\end{table}
\begin{longtable}{lrr}
\caption{Распределение документов и текста по источникам.}\\
\toprule Источник & Документов & Текст, MiB\\\midrule\endfirsthead
\toprule Источник & Документов & Текст, MiB\\\midrule\endhead
clang & 144 & 4,56\\
freebsd-docs & 183 & 12,60\\
linux-kernel & 3\,516 & 37,41\\
cppreference & 6\,838 & 20,41\\
gcc & 314 & 2,17\\
llvm-libc & 133 & 0,45\\
llvm-libcxx & 38 & 0,44\\
lld & 14 & 0,06\\
lldb & 9\,598 & 85,21\\
llvm & 314 & 6,49\\
man7 & 5\,473 & 22,59\\
binutils & 4 & 0,02\\
gdb & 1 & 0,02\\
glibc & 306 & 1,20\\
codereview.stackexchange.com & 33\,980 & 461,43\\
softwareengineering.stackexchange.com & 21\,456 & 240,22\\
unix.stackexchange.com & 18\,921 & 184,43\\
dba.stackexchange.com & 9\,995 & 101,76\\
serverfault.com & 19\,463 & 192,03\\
askubuntu.com & 22\,181 & 242,55\\
superuser.com & 26\,517 & 242,14\\
stackoverflow.com & 1\,500\,582 & 14579,69\\
\bottomrule
\end{longtable}


Stack Overflow составляет около 89,3\% документов. Поэтому даже большое
число источников не означает равномерное покрытие тем. Суммарная статистика
также не означает, что каждый документ содержит несколько тысяч слов:
среди справочных страниц есть короткие описания API.

Для внешнего поиска подходят запросы
\code{site:docs.kernel.org memory allocation},
\code{site:man7.org pthread_create},
\code{site:en.cppreference.com std::vector} и поиск
\code{[c++] memory allocation} по Stack Overflow.
Они проверяют соответственно тему, точное имя функции и технический
идентификатор. Ограничения такого сравнения: внешняя система индексирует
другую версию сайтов, может включать отсутствующие в архиве страницы и
ранжирует их по собственным сигналам. Сохранённого снимка внешней выдачи
нет, поэтому этим примерам не приписываются измеренные Precision или NDCG.

\conclusions
Получен тематический корпус, который можно обрабатывать потоком и который
превышает миллион документов. Объём извлечённого текста составляет около
16,05 GiB. Основные ограничения - преобладание Stack Overflow,
неодинаковая длина документов и возможные содержательные дубликаты.
Улучшения: дедупликация по нормализованному содержимому, стратифицированная
оценка по источникам и отдельный контроль качества извлечения кода.

\lab{2}{Поисковый робот}
\task
Разработать робота, принимающего единственный аргумент - путь к YAML-конфигу
с секциями \code{db} и \code{logic}. Сохранять нормализованный URL,
исходный HTML, источник и время загрузки Unix timestamp. Обеспечить
продолжение после остановки, периодическую проверку документов и обновление
содержимого только при изменении.

\method
Робот реализован на Python с \code{asyncio}, HTTPX и MongoDB.
Конфигурация задаёт соединение с базой, задержки, тайм-аут, период
повторной проверки и источники с начальными URL. Внутри источника
действуют разрешённые и исключённые префиксы пути. Навигационная страница
может служить отправной точкой, не попадая в итоговую коллекцию документов.

Нормализация URL приводит домен к нижнему регистру и IDNA, удаляет
стандартный порт и фрагмент, задаёт пустому пути значение \code{/}.
Параметры запроса сохраняются. Это устраняет часть повторов адресов,
но не решает общую задачу эквивалентности URL с разными query-параметрами.

Очередь хранится в MongoDB. Задания переходят между состояниями
\code{pending}, \code{processing}, \code{done}, \code{failed},
\code{blocked}. Атомарный захват задания предотвращает его одновременную
обработку. При новом запуске незавершённые задания возвращаются в очередь,
а завершённые перепроверяются после истечения заданного интервала.
Документ, найденные ссылки и завершение задания фиксируются транзакцией;
поэтому MongoDB запускается как одноузловой replica set.

Разные хосты обходятся параллельно, запросы к одному хосту - последовательно.
Задержка отсчитывается между началами запросов и учитывает ограничения
источника и \code{robots.txt}. Наибольшая задержка применяется ко всем
источникам одного хоста. Перенаправления обрабатываются явно, а новый адрес
повторно проверяется на допустимость обхода.

Условные HTTP-запросы используют ETag и Last-Modified. Ответ 304 не
перезаписывает содержимое. При ответе с телом сравнивается SHA-256:
одинаковый HTML повторно не сохраняется. Время последней проверки и время
изменения содержимого имеют разный смысл. Ошибка загрузки не должна удалять
уже сохранённую версию документа.

\journal
Логика разделена между конфигурацией, HTTP-клиентом, правилами источников,
хранилищем и движком обхода. После реализации последовательной загрузки
добавлены очередь в базе, восстановление и повторная проверка. Затем
введены ограничения \code{robots.txt}, общий регулятор задержек хоста
и конкурентный обход разных хостов. Сетевое ожидание вынесено из блокировки
операций БД, чтобы медленный сайт не останавливал остальные источники.

Аварийное завершение между отправкой HTTP-запроса и транзакцией может
привести к повторной загрузке страницы. Полностью исключить её нельзя без
транзакции с удалённым HTTP-сервером, но можно гарантировать отсутствие
потери задания и дублирования документа в базе.

\testing
\begin{longtable}{p{0.32\textwidth}p{0.60\textwidth}}
\toprule Сценарий & Ожидаемое поведение\\\midrule
Нормализация URL & Фрагмент удалён, query сохранён, стандартный порт убран.\\
Остановка и перезапуск & Задание processing снова доступно для обработки.\\
Повторный HTML, 304 & Тело документа не перезаписано.\\
Изменённый HTML & Новое содержимое и хеш сохранены.\\
Ошибки HTTP и тайм-аут & Старая версия сохранена, ошибка отмечена в очереди.\\
Запрет robots.txt & Документ не загружается в обход запрета.\\
Редирект за границы источника & Целевой адрес проверяется до загрузки.\\
Два источника одного хоста & Соблюдается общая задержка запросов.\\
Разные хосты & Сетевые ожидания могут перекрываться.\\\bottomrule
\end{longtable}
Эти группы представлены в \code{2_crawler/tests}: отдельные модули
проверяют очередь, хранение, повторный обход, robots, редиректы и задержки.
Автотесты запускаются командой \code{just test}. Они проверяют
правила и переходы состояния робота. Для проверки транзакций
с настоящей MongoDB нужен отдельный интеграционный запуск.

\results
Реализованы требуемый интерфейс конфигурации, возобновляемая очередь и
проверка изменений. Минимальный запуск:
\begin{lstlisting}
cd 2_crawler
docker compose up -d --wait
just run config.yaml
just test
\end{lstlisting}
Штатный пример задаёт задержку 2 секунды глобально и 3 секунды для Linux
Kernel, тайм-аут 30 секунд и повторную проверку через 86\,400 секунд.
Следовательно, задержка источника ограничивает частоту начала запросов
примерно одним запросом за три секунды без учёта более строгих правил сайта.
Это ограничение конфигурации, а не измеренная скорость выгрузки.

\conclusions
Сохранение очереди отдельно от процесса позволяет продолжать обход после
сбоя. Условные запросы сокращают передачу неизменившихся страниц,
а транзакции согласуют документ и очередь. Ограничения решения -
одно задание хоста за раз, отсутствие распределённой координации нескольких
роботов и необходимость MongoDB replica set. Большая часть итогового корпуса
получена импортом архивов, поэтому её объём нельзя считать результатом
HTTP-производительности робота.

\lab{3}{Токенизация}
\task
Разработать правила выделения токенов и реализацию для дальнейшей
индексации. Описать достоинства, недостатки и неудачные случаи. Измерить
число токенов, среднюю длину, время обработки и скорость на килобайт;
обсудить зависимость от размера данных и способы ускорения.

\method
Токенизатор выделяет последовательности букв Unicode, десятичных цифр
и подчёркиваний. Комбинируемые знаки допускаются после начала сегмента.
Последовательность из одних подчёркиваний не выдаётся. Регистр понижается
по таблице Unicode 15.0; полное case folding и NFC-нормализация не
выполняются. Порядок и повторы сохраняются, стоп-слова не удаляются.

Для технического корпуса введены два дополнения. Разделитель
\code{::} сохраняется между корректными сегментами, поэтому
\code{std::vector::iterator} остаётся одним термином. Обозначения
\code{C++}, \code{G++}, \code{Clang++}, в том числе с номером версии,
также сохраняются. Обычное увеличение \code{value++} таким правилом
не защищается. Это компромисс между естественным текстом и исходным кодом,
а не полноценный синтаксический анализ C++.

\begin{table}[H]\centering
\begin{tabularx}{\textwidth}{lX}\toprule
Вход & Результат\\\midrule
\code{Memory allocation} & \code{memory}, \code{allocation}\\
\code{pthread_create} & \code{pthread_create}\\
\code{std::vector<int>} & \code{std::vector}, \code{int}\\
\code{C++20 value++} & \code{c++20}, \code{value}\\
\code{memory-mapped} & \code{memory}, \code{mapped}\\
\code{can't -O2 6.12} & \code{can}, \code{t}, \code{o2}, \code{6}, \code{12}\\
\code{MemoryBarrier} & \code{memorybarrier}\\\bottomrule
\end{tabularx}
\caption{Примеры правил и их ограничений.}
\end{table}

Строгий UTF-8-декодер отвергает обрезанные последовательности,
недопустимые кодовые точки и ошибочные байты. Длина токена считается
в кодовых точках, включая сохранённые символы \code{_}, \code{:},
\code{+}. Это не длина UTF-8 в байтах и не число визуальных графем.
Среднее вычисляется по всем вхождениям:
\[
\overline{\ell}_{token}=\frac{\sum_{i=1}^{T}\ell(t_i)}{T}.
\]
Библиотечная функция вызывает callback для каждого токена. Самостоятельные
CLI и обработчик \code{corpus.bin} используют один и тот же код.

\journal
Базовые правила для текста дополнены сохранением имён API и обозначений
языка. Затем добавлены таблицы Unicode, строгая проверка входа и
статистический режим без печати каждого токена. Измеритель потока отдельно
считает время внутри вызовов токенизатора и полное время чтения корпуса.
При включённых частотах callback обновляет словарь, поэтому такой замер
нельзя напрямую сравнивать с запуском без сбора частот.

\testing
Табличные C++-тесты покрывают пустой ввод, пробелы, пунктуацию, повторы,
стоп-слова, идентификаторы, пространства имён, версии, пути, числа,
кириллицу, диакритику, комбинируемые знаки и символы вне BMP.
Проверяются не только строки, но и длины. Отдельные случаи неверного
UTF-8 должны завершаться исключением. Python-тесты проверяют извлечение,
профили сайтов, экспорт и чтение корпуса. Таким образом, токенизация и
качество предварительного извлечения проверяются раздельно.

\results
\begin{table}[H]\centering
\begin{tabular}{lrr}\toprule
Показатель & Префикс 10\,000 & Весь корпус\\\midrule
Документов & 10\,000 & 1\,679\,971\\
Текст, байт & 76\,863\,027 & 17\,236\,388\,223\\
Токенов & 10\,683\,808 & 2\,483\,697\,508\\
Средняя длина, кодовых точек & 5,5300 & 4,8031\\
Время токенизации, с & 3,073 & 401,114\\
Полное время, с & 3,392 & 434,245\\
Скорость, MiB/с & 23,853 & 40,981\\
\bottomrule\end{tabular}
\caption{Повторные измерения токенизации. В префиксе включён сбор частот; на полном корпусе он отключён.}
\end{table}
Для полного корпуса время на KiB равно 23,830 мкс. Время
токенизации не включает чтение и расчёт отпечатка, полное время включает их.
Поле \code{terms=0} полного запуска означает отключённый сбор словаря,
а не отсутствие уникальных термов.

В основном индексе сохранено 2\,483\,697\,508 позиций токенов:
эта величина относится ко всему корпусу и совпадает у булевой и
координатной индексации. Она больше числа слов экспортёра, поскольку,
например, дефис, точка и некоторые элементы программного текста
создают дополнительные границы.

По длине входа основной проход линейный. Поиск свойств символа использует
бинарный поиск по фиксированным таблицам Unicode; его стоимость ограничена
размером таблиц, а не корпуса. Для проверки масштабирования можно
использовать несколько префиксов одного и того же потока, сохраняя
одинаковые настройки сбора частот. Скорость первого префикса не следует
считать скоростью случайного среднего документа: в начале находятся
преимущественно страницы документации.

\conclusions
Получен общий токенизатор для корпуса и запросов. Сохранение технических
обозначений уменьшает потерю смысла, но \code{std::vector} и
\code{vector} остаются разными словами, а апострофы, пути и ключи
командной строки обрабатываются грубо. Возможны дополнительные формы
идентификатора, разбиение CamelCase и NFC-нормализация; они потребуют
согласованного изменения индекса и запросов. Ускорить обработку можно
быстрым путём для ASCII и параллельной обработкой независимых документов.

\lab{4}{Лемматизация и стемминг}
\task
Добавить поиск без учёта словоформ; сравнить поведение с точным поиском,
рассмотреть ухудшившиеся запросы и возможный бонус за точное совпадение.
Описать выбранную нормализацию и оценку качества после её включения.

\method
Реализованы три режима: \code{exact}, \code{stem} и \code{lemma}.
Первый оставляет токены после понижения регистра. Второй использует
самостоятельную реализацию английского стеммера Портера. Третий применяет
словарь WordNet 3.0 и собственные правила получения кандидатов.
Одинаковая функция нормализации применяется при индексации и поиске.

Портер последовательно преобразует суффиксы, учитывая наличие гласной,
меру чередования гласных и согласных, удвоенные согласные и окончания CVC.
Результат не обязан быть словом языка: \code{connections} даёт
\code{connect}. Слишком короткие слова и токены с символами вне
\code{a--z} не меняются. Поэтому \code{pthread_create},
\code{std::vector}, числа и русские токены не повреждаются.

Словарный метод сначала проверяет исключения и наличие исходной формы,
затем строит кандидатов по окончаниям для существительных, глаголов,
прилагательных и наречий. Замена принимается только при наличии кандидата
в словаре с нужным признаком части речи. Так обрабатываются нерегулярные
формы \code{children -> child}, \code{mice -> mouse},
\code{went -> go}. Контекст предложения и снятие омонимии не используются.
При неоднозначности важен фиксированный порядок правил.

WordNet используется как набор данных, а не как готовая поисковая
библиотека. Скрипт \code{build_lexicon.py} упаковывает словарные записи
и исключения; CMake встраивает массив в исполняемый файл.
Собственный бинарный поиск работает по упакованному словарю.
Идентификатор словаря записывается в координатный индекс, чтобы исключить
его чтение с несовместимой нормализацией.

\journal
Сначала реализована защита идентификаторов и каскад Портера, затем
добавлены словарь, исключения и морфологические правила WordNet.
Для сравнения сохранён режим без морфологии. В булевом индексаторе
предусмотрен дополнительный exact-индекс: сначала выводятся документы
с точной формой, затем остальные совпадения по основе. В штатной
конфигурации дополнительный индекс отключён. Лабораторные 7--8 используют
стемминг, 10--13 и 17 - WordNet.

\testing
Проверяются классические примеры шагов Портера, нерегулярные формы,
слова без преобразований, пустая строка и сохранение идентификаторов.
Особый случай - увеличение длины при \code{mice -> mouse}:
при достаточном буфере преобразование проходит, при недостаточном
должна возникать ошибка. Для качества нужно сравнить одинаковые запросы
и корпус в трёх режимах, разметить объединённый пул и оценить метрики
из лабораторной №6. Совпадение основы само по себе не доказывает
релевантность найденного документа.

\results
\begin{table}[H]\centering
\begin{tabular}{lll}\toprule
Вход & Портер & WordNet\\\midrule
\code{connections} & \code{connect} & \code{connection}\\
\code{allocations} & \code{alloc} & \code{allocation}\\
\code{children} & \code{children} & \code{child}\\
\code{running} & \code{run} & \code{run}\\
\code{pthread_create} & \code{pthread_create} & \code{pthread_create}\\\bottomrule
\end{tabular}
\caption{Ожидаемые преобразования, задаваемые правилами и тестами.}
\end{table}
В полном булевом индексе со стеммингом 19\,787\,654 терма, в
координатном индексе с WordNet - 20\,430\,767. Разница составляет
643\,113 термов (около 3,15\% от словаря WordNet), но это сравнение
двух способов нормализации. Оно не измеряет сокращение относительно
отдельно построенного exact-индекса.

В выдаче по полному корпусу запросы \code{memory allocation} и
\code{memory allocations} дают одинаковые первые результаты TF-IDF
и по 282\,185 кандидатов: формы объединяются. Для запроса
\code{epolll} исправление опечатки не выполняется и результатов нет.
Оценка качества системы с WordNet приведена в лабораторной №6.
Изолированное влияние морфологии требует отдельного сравнения
exact и lemma при одинаковом отборе и ранжировании кандидатов.
Сравнение строгого поиска и TF-IDF такой эксперимент не заменяет.

\conclusions
Нормализация объединяет словоформы, но может объединять разные смыслы
и оставлять часть родственных слов раздельными. Словарь полезен для
нерегулярных форм, однако не решает проблему контекста и опечаток.
Для технических терминов сохранение точной формы и аккуратное расширение
запроса могут быть полезнее агрессивного сведения к основе.
Изолированное сравнение режимов следует проводить на общем пуле
с оценками релевантности для каждого варианта.

\lab{5}{Закон Ципфа}
\task
Построить распределение частот терминов в логарифмической шкале,
наложить закон Ципфа и объяснить расхождения. Необязательная часть -
подбор параметров закона Мандельброта.

\method
Для каждого уникального токена считаются число вхождений \(cf\) и число
документов \(df\). По убыванию \(cf\) назначаются ранги \(r=1,\ldots,V\);
равные частоты упорядочиваются лексикографически для воспроизводимости.
В этой работе морфологическая нормализация не включается, чтобы
отделить свойства токенизации от стемминга.

Модель Ципфа имеет вид \(f(r)=C/r\). Константа выбирается из сохранения
общей массы частот:
\[
H_V=\sum_{r=1}^{V}\frac1r,\qquad C=\frac{T}{H_V},
\qquad\sum_{r=1}^{V}\frac{C}{r}=T.
\]
Это не подгонка по самому частому слову. Дополнительно методом наименьших
квадратов оценивается наклон зависимости \(\ln f\) от \(\ln r\)
по всем термам. В идеальной модели он равен \(-1\).
Для отображения можно проредить ранги; расчёт регрессии выполняется
до прореживания. Гапакс - термин с одним вхождением во всём наборе.

\journal
Потоковый токенизатор дополнен частотами и множеством термов текущего
документа, чтобы \(df\) увеличивалась лишь один раз на документ.
Скрипт \code{5_zipf/tools/zipf.py} проверяет уникальность строк,
\(cf\ge df\ge1\), сортирует термы, считает модель и регрессию.
Статистика вычислена для первых 10\,000 документов
корпуса. Поскольку регрессия и сортировка требуют полного словаря
исследуемого набора, для этого опыта выбран ограниченный префикс.
Полученные частоты и наклон характеризуют именно эти 10\,000 документов.

\testing
Для искусственного словаря проверяются ранги, сумма частот и значение
константы. Пустой словарь, дубликат терма и \(cf<df\) должны отвергаться.
При совпадении частот порядок должен оставаться стабильным. Для
одного терма наклон условно равен нулю, поскольку дисперсия логарифмов
ранга равна нулю. При прореживании необходимо сохранять первый и
последний ранги. Перечисленные проверки задают план тестирования
частотного анализа и построения графика.

\results
В эксперименте: 10\,000 документов, 10\,683\,808 вхождений,
401\,968 уникальных токенов, 266\,496 гапакс
(66.30\% словаря).
Константа $C=792488.31$, наклон регрессии $-1.0655$.
Это характеристики выборки без стемминга, а не всего корпуса.
Самые частые токены: \code{the} (544\,244), \code{to} (230\,087), \code{a} (192\,740), \code{is} (175\,251), \code{of} (170\,951).

\begin{figure}[H]\centering
\begin{tikzpicture}
\begin{loglogaxis}[width=0.95\textwidth,height=7cm,
xlabel={Ранг термина},ylabel={Число вхождений},grid=both,
legend style={at={(0.02,0.02)},anchor=south west,font=\small}]
\addplot[blue,thick] table[x=rank,y=count]{figures/zipf.dat};
\addlegendentry{Наблюдаемые частоты}
\addplot[red,dashed,thick] table[x=rank,y=zipf]{figures/zipf.dat};
\addlegendentry{Модель $C/r$}
\end{loglogaxis}
\end{tikzpicture}
\caption{Распределение для первых 10\,000 документов. Все точки получены из фактических частот этой выборки.}
\end{figure}
Частые слова образуют начало распределения, технические имена и
уникальные обозначения - длинный хвост. При малых целых частотах
видны ступени. Отклонение от прямой объясняется смесью прозы и кода,
ограниченностью выборки, повторяемостью документации и отсутствием
фильтра стоп-слов. Линейная регрессия в логарифмах является описанием
формы кривой, а не статистическим доказательством степенного закона.

\conclusions
Распределение имеет длинный хвост редких терминов, поэтому размер словаря
может расти значительно быстрее, чем подсказывает число частых слов.
Модель Ципфа полезна как приближение и объясняет неравномерную длину
списков постингов. Опыт ограничен префиксом корпуса с документацией;
для вывода о всём корпусе нужен полный пересчёт или стратифицированная
выборка обсуждений и документации. Закон Мандельброта не подбирался.

\lab{6}{Оценка качества поиска}
\task
Оценить собственный поиск и две альтернативы на 30 запросах.
Вычислить P, DCG, NDCG и ERR на глубинах 1, 3 и 5; разобрать сильные
и слабые стороны систем и возможные улучшения.

\method
Собраны 30 запросов о памяти, синхронизации, процессах, файловых
дескрипторах и API. Набор включает короткие имена \code{fork},
\code{malloc}, \code{pthread_create}, словоформы
\code{memory allocation / memory allocations}, более широкие
\code{resource limits} и опечатку \code{epolll}. Он отражает несколько
типов сложности, но не является случайной выборкой пользовательского лога.

Сравниваются собственный строгий поиск, собственный TF-IDF и две
альтернативы: SQLite FTS5 и Tantivy. Готовые библиотеки используются
только для внешних базовых систем, собственный индекс ими не строится.
Выдачи сохраняются в TSV: система, запрос, ранг, URL. Релевантность
отдельно задаётся для пары «запрос, URL» по шкале 0--3: нерелевантный,
слабосвязанный, полезный, непосредственно отвечающий на запрос.
Для Precision релевантными считаются оценки не ниже 2.

Пусть \(g_i\) - оценка документа на позиции \(i\). Тогда:
\begin{align*}
P@k&=\frac1k\sum_{i=1}^{k}[g_i\ge2],\\
DCG@k&=\sum_{i=1}^{k}\frac{2^{g_i}-1}{\log_2(i+1)},\\
NDCG@k&=\frac{DCG@k}{IDCG@k},\\
ERR@k&=\sum_{i=1}^{k}\frac{R_i}{i}\prod_{j=1}^{i-1}(1-R_j),
\qquad R_i=\frac{2^{g_i}-1}{8}.
\end{align*}
Идеальный DCG получается сортировкой всех известных оценок запроса.
В оценочный набор входят разметка объединённого пула и дополнительные
эталонные документы. Поэтому IDCG не ограничен результатами одной системы.
При \(IDCG=0\) NDCG определён нулём. Короткая выдача дополняется
нулевыми оценками до глубины отсечения; её Precision всё равно делится
на \(k\). Неоценённый выданный документ принципиально отличается
от отсутствующего результата: программа завершает расчёт ошибкой.
Итоговые метрики усредняются по запросам, каждый запрос имеет одинаковый вес.

\journal
Сначала реализованы формулы и ручные контрольные примеры, затем
чтение TSV и проверка входных данных. Для сравнения построены индексы
SQLite и Tantivy на тех же 1\,679\,971 документах. Для каждого запроса
объединены первые пять результатов всех систем, повторные URL удалены.
Все 426 уникальных пар «запрос, документ» явно оценены с обоснованиями.
Оценивание учитывает не только заголовок, но и текст документа;
для неоднозначных длинных страниц дополнительно проверяются фрагменты
об искомом API. Контрольные суммы входных данных защищают от смешения
выдач и оценок разных экспериментов.

\testing
Ручной пример \(g=(3,0,2)\), идеальный список \((3,2)\), \(k=3\):
\[
P@3=\frac23,\quad DCG@3=8{,}5,\quad
NDCG@3=\frac{8{,}5}{7+3/\log_2 3},\quad
ERR@3=0{,}875+\frac{0{,}125\cdot0{,}375}{3}.
\]
Дополнительно проверяются пустая выдача, короткая выдача и NUL
в строке входа. Три модульных теста пройдены.
Дополнительные семь Python-тестов проверяют полноту пула,
отклонение пропущенных и повторных оценок, допустимость шкалы,
наличие обоснования, покрытие запросов и контрольные суммы.
Повторный расчёт метрик выполнен без чтения корпуса и индексов.

\results
\begin{table}[H]\centering
\begin{tabular}{lrrr}\toprule
Система & Результатов & Оценено & Покрытие, \%\\\midrule
Собственный булев поиск & 145 & 145 & 100,00\\
Собственный TF-IDF & 145 & 145 & 100,00\\
SQLite FTS5 & 145 & 145 & 100,00\\
Tantivy & 145 & 145 & 100,00\\\bottomrule
\end{tabular}
\caption{Покрытие оценками первых пяти результатов 30 запросов.}
\end{table}
145 вместо 150 результатов связано с пустой выдачей запроса
\code{epolll}. Он включён в усреднение и имеет нулевые метрики.
\begin{longtable}{llrrrr}
\caption{Средние метрики по 30 запросам.}\\
\toprule Система & Глубина & P & DCG & NDCG & ERR\\\midrule\endfirsthead
\toprule Система & Глубина & P & DCG & NDCG & ERR\\\midrule\endhead
Булев поиск & @1 & 0,1000 & 0,8333 & 0,1190 & 0,1042\\
Булев поиск & @3 & 0,1667 & 2,0484 & 0,1469 & 0,1942\\
Булев поиск & @5 & 0,1933 & 2,9820 & 0,1744 & 0,2283\\
\midrule
TF-IDF & @1 & 0,4333 & 2,9333 & 0,4190 & 0,3667\\
TF-IDF & @3 & 0,4333 & 5,2305 & 0,3557 & 0,4423\\
TF-IDF & @5 & 0,4933 & 7,8437 & 0,4219 & 0,4742\\
\midrule
SQLite FTS5 & @1 & 0,3667 & 2,1667 & 0,3095 & 0,2708\\
SQLite FTS5 & @3 & 0,4667 & 5,1214 & 0,3578 & 0,4257\\
SQLite FTS5 & @5 & 0,4067 & 6,5297 & 0,3745 & 0,4448\\
\midrule
Tantivy & @1 & 0,4000 & 2,4000 & 0,3429 & 0,3000\\
Tantivy & @3 & 0,4667 & 5,1563 & 0,3608 & 0,4421\\
Tantivy & @5 & 0,4267 & 6,9174 & 0,3937 & 0,4716\\
\bottomrule
\end{longtable}

На глубине 5 собственный TF-IDF имеет наибольшие значения всех
четырёх метрик среди сравниваемых систем. Его P@5 равна 0,4933
против 0,1933 строгого поиска, NDCG@5 равна 0,4219 против 0,1744.
При этом на глубине 3 Precision у SQLite FTS5 и Tantivy выше:
0,4667 против 0,4333. Поэтому преимущество зависит от глубины и метрики.

Разбор отдельных запросов показывает разные типы ошибок.
На \code{malloc} собственный TF-IDF достигает NDCG@5 = 1,0000,
а на \code{reference counting} получает 0,0000. В последнем случае
SQLite FTS5 даёт NDCG@5 = 0,4524. Общие слова и широкое объединение
кандидатов могут вытеснять документы, отвечающие именно на задачу
пользователя. Для исправления полезны учёт словосочетаний,
близости слов и специфики справочных страниц.

Оценки выставлены ассистентом без независимых человеческих оценщиков.
Кроме того, пул ограничен первыми пятью результатами сравниваемых
систем. Полное покрытие этого пула не означает полной разметки корпуса.

\conclusions
Выполнено сравнение четырёх систем на общем размеченном наборе.
TF-IDF существенно улучшил качество относительно порядка выдачи
строгого поиска, однако не оказался лучшим для каждого запроса
и каждой глубины. Чтобы повысить надёжность выводов, нужны
несколько оценщиков, расширение пула и более разнообразный набор запросов.

\lab{7}{Булев индекс}
\task
Построить собственный бинарный обратный индекс и прямой индекс заголовков
и URL. Описать побайтовый формат, внутреннее представление, сортировку,
количество и среднюю длину термов, ограничения скорости и масштабирования.

\method
Корпус читается последовательно. Каждому документу назначается
\(docID\) по порядку в потоке. Токенизация и стемминг создают записи
«терм, документ, частота». Буфер ограничен параметром \code{memory_mib}.
После заполнения он сортируется по терму и \(docID\), одинаковые пары
объединяются суммированием частот, результат записывается во временный
run-файл. Отсортированные файлы попарно сливаются до единого потока.
Это вариант блочной индексации с внешней сортировкой.

Для массивов, сортировки и операций со списками используются собственные
структуры и функции. Файл записывается явно по полям, а не выгрузкой
структуры C++: выравнивание компилятора не определяет формат.
Все многобайтные целые имеют порядок little-endian.

\begin{longtable}{rlp{0.53\textwidth}}
\caption{Заголовок булевого индекса IRIDX001, 96 байт.}\\
\toprule Смещение & Размер & Поле\\\midrule\endfirsthead
\toprule Смещение & Размер & Поле\\\midrule\endhead
0 & 8 & Сигнатура \code{IRIDX001}\\
8 & 4 & Версия формата: 1\\
12 & 4 & Флаги нормализации\\
16 & 8 & Число документов\\
24 & 8 & Число термов\\
32 & 8 & Число токенов\\
40 & 8 & Объём входного текста\\
48 & 8 & Начало словаря\\
56 & 8 & Начало строк термов\\
64 & 8 & Начало постингов\\
72 & 8 & Начало смещений документов\\
80 & 8 & Начало метаданных\\
88 & 8 & Размер файла\\\bottomrule
\end{longtable}
Запись терма занимает 32 байта: смещение текста (8), смещение списка
постингов (8), суммарная частота (8), длина UTF-8 в байтах (4),
число документов (4). Смещения терма отсчитываются от соответствующих
секций. Словарь отсортирован лексикографически; поиск выполняется
бинарным поиском. Список постингов содержит возрастающие \code{uint32}
идентификаторы, по 4 байта на документ. Частота хранится агрегатом
в словаре, отдельной \(tf\) в булевом постинге нет.

Прямой индекс содержит 8-байтовые смещения записей документов.
В каждой записи идут две 4-байтовые длины, затем URL и заголовок.
Отделение прямого индекса позволяет извлекать метаданные только для
показываемых результатов. Флаги и версия создают основу для расширения,
но произвольный новый формат старый читатель не поддержит автоматически.

\journal
Построение разделено на накопление блока, сброс, слияние и сборку секций.
Внутри блока одинаковые записи схлопываются до записи на диск.
Для защиты данных добавлены проверки размеров, порядка постингов,
диапазона идентификаторов и границ секций. Корпус читается потоком,
поэтому его полный размер не определяет объём оперативной памяти.
При этом лимит блока не является строгим лимитом RSS всего процесса:
дополнительная память нужна строкам, буферам и служебным структурам.

\testing
Тест \code{ExternalSortAndMetadata} строит один корпус с буфером
128 байт и 1 MiB и сравнивает файлы побайтово. Проверяются суммарная
частота, число документов терма, отсутствие неизвестного слова и
заголовок Unicode. Второй тест проверяет NUL и обрезанные данные.
Оба теста пройдены. Эти тесты проверяют независимость
результата от разбиения на блоки, а не скорость на полном корпусе.

\results
\begin{table}[H]\centering
\begin{tabular}{lr}\toprule
Показатель & Значение\\\midrule
Документов & 1\,679\,971\\
Токенов & 2\,483\,697\,508\\
Термов после стемминга & 19\,787\,654\\
Средняя длина терма, кодовых точек & 13,2517\\
Размер булевого индекса, байт & 3\,618\,409\,348\\\bottomrule
\end{tabular}
\caption{Статистика готового булевого индекса, повторно прочитанная из файла.}
\end{table}
Средняя длина уникального терма существенно больше средней длины
вхождения токена: частые короткие служебные слова в первой величине
имеют вес один, а длинные редкие идентификаторы образуют большую часть
словаря. Стемминг укорачивает некоторые слова, но не устраняет этот эффект.

Время полного построения булевого индекса отдельно не измерялось.
Поэтому его скорость не выводится из времени сборки координатного индекса.
Параметр 64 MiB в конфигурации - размер рабочего бюджета блока,
а не измеренное пиковое потребление памяти. При росте корпуса в 10,
100 и 1000 раз растут число временных блоков и объём слияния.
Дисковая пропускная способность и число проходов становятся ограничением;
хранение \(docID\) в 32 битах дополнительно ограничивает число документов.

\conclusions
Создан переносимый бинарный формат с прямым и обратным индексом.
Внешняя сортировка позволяет строить его без загрузки всех документов
в память. Слабое место - повторное чтение временных файлов при
попарном слиянии. Многопутевое слияние и уменьшение числа проходов
могут снизить дисковую нагрузку. Для полноценной оценки скорости
необходимо сохранять журнал построения вместе со статистикой машины.

\lab{8}{Булев поиск}
\task
Реализовать запросы с И (пробел или \code{&&}), ИЛИ (\code{||}),
НЕ (\code{!}) и скобками; CLI с запросами по строкам stdin и веб-форму
с заголовками, URL и постраничной выдачей. Проверить корректность
и обсудить сложные запросы и скорость.

\method
Рекурсивный парсер строит дерево запроса. Приоритеты:
\code{!}, затем \code{&&}, затем \code{||}. Скобки задают другой
порядок, соседние атомы в строгом режиме соединяются операцией И.
Пустой запрос даёт пустое множество. Одиночные \code{&}, незакрытые
скобки и незавершённая операция считаются синтаксической ошибкой.

Терм нормализуется так же, как при индексации. Его постинг-лист читается
из файла; результат поддерева - отсортированное множество \(docID\).
Пересечение и объединение выполняются двумя указателями за
\(O(a+b)\), отрицание - дополнением относительно диапазона всех
документов \([0,N)\). Поэтому пустой документ может входить в ответ
на \code{!memory}. Повторы идентификаторов в выдаче исключаются.

CLI загружает индекс и принимает несколько запросов за запуск.
Ответ JSON содержит общее число совпадений, смещение, время,
число прочитанных постингов и результаты. Веб-обвязка запускает
локальный поиск и экранирует заголовки и ссылки при формировании HTML.
Параметры \code{search_limit=50} и \code{search_offset} реализуют
постраничное извлечение. Веб-сервер запускается отдельно от индексатора.

\journal
Разбор запроса отделён от операций над множествами. Это позволило
позднее сохранить язык булевых операций при переходе к координатному
индексу. Добавлены обработка переменного количества пробелов,
двойное отрицание и диагностика некорректного синтаксиса.
Нормализация запросов перенесена в общий конвейер токенизации,
чтобы запрос не расходился с представлением документа.

\testing
На искусственном корпусе из четырёх документов проверяются все
операции, приоритеты, пустой документ, неизвестный терм и пустой ввод.
Для \code{memory || connections} ожидается объединение без повторов;
для \code{connections !mutex} - исключение документа с mutex.
Дерево \code{alpha || beta && !delta} должно иметь корневую операцию
ИЛИ и правую ветвь И. Два модульных теста пройдены.

Для веб-интерфейса дополнительно требуются ручные проверки пустой
выдачи, перехода к следующей странице, спецсимволов в заголовке
и ошибки запроса. Эти проверки не заменяются C++-тестами парсера.

\results
Примеры запросов для демонстрации:
\begin{lstlisting}
memory allocation
memory && (allocation || mapping)
(mutex || spinlock) && !futex
!memory
\end{lstlisting}
\begin{lstlisting}
cd 8_boolean
just build
just search < queries.txt
just web
\end{lstlisting}
Готовый индекс уже существует; повторное \code{build-index} для поиска
не требуется. Время работы зависит от количества прочитанных
постингов и размеров промежуточных множеств. Замеры координатного
движка из следующей работы не характеризуют скорость этой реализации.

Дорогие запросы - объединения частых слов, чистое отрицание и широкие
подвыражения, создающие промежуточные списки почти из всех документов.
Выдача только первых 50 заголовков не устраняет стоимость полного
вычисления множества совпадений в текущем алгоритме.

\conclusions
Язык запросов позволяет точно задавать условия присутствия слов,
но порядок выдачи по \(docID\) слабо связан с полезностью документа.
Возможные улучшения - выполнение И начиная с самого короткого
списка, курсоры вместо материализации промежуточных множеств,
прыжки и ранжирование. Эти направления реализуются в следующих работах.

\lab{10}{Цитатный поиск и координатный индекс}
\task
Расширить булев поиск точными цитатами и упорядоченными цитатами
\code{"..." / N}, где расстояние между первым и последним словом
не превышает \(N\). Описать координатный формат, размер, число
позиций, средние количества позиций, скорость и примеры запросов.

\method
Для каждой пары «терм, документ» хранится возрастающий список позиций
вхождений. Позиция - номер токена после токенизации, а не смещение
символа. Стоп-слова не удаляются, поэтому координаты соседних слов
сохраняют исходное расстояние. Нормализация меняет текст терма,
но не количество его вхождений.

Кандидаты цитаты должны содержать все её термы. В качестве ведущего
берётся терм с минимальным \(df\); курсоры остальных списков продвигаются
к текущему документу. Для точной цитаты длины \(m\) проверяются
позиции \(p,p+1,\ldots,p+m-1\). Для цитаты с расстоянием
выбираются строго возрастающие позиции \(p_1<\cdots<p_m\),
удовлетворяющие \(p_m-p_1\le N\). Повтор слова в цитате требует
разных последовательных вхождений. Порядок слов не переставляется.

Текущий формат - \code{IRPOS003}. Общий заголовок занимает 144 байта.
Первые 96 байт имеют назначение, аналогичное булевому формату,
но сигнатура и версия другие. Продолжение:
\begin{table}[H]\centering
\begin{tabular}{rll}\toprule
Смещение & Размер & Поле\\\midrule
96 & 8 & Секция длин документов\\
104 & 8 & Секция прыжков\\
112 & 8 & Секция норм векторов\\
120 & 8 & Отпечаток корпуса\\
128 & 4 & Шаг прыжков\\
132 & 4 & Идентификатор словаря WordNet\\
136 & 8 & Размер заголовка: 144\\\bottomrule
\end{tabular}
\caption{Расширение заголовка координатного индекса.}
\end{table}
Терм занимает 56 байт: пять 8-байтовых полей (смещение текста,
смещение постингов, длина постингов, смещение прыжков, суммарная
частота) и четыре 4-байтовых поля (длина текста, \(df\), число
блоков, резерв). Постинг включает \(docID\), число вхождений и
позиции. В raw-режиме каждое число имеет размер 4 байта;
VB-режим описан в следующей работе.

Прямой индекс дополнен длинами документов и нормами для TF-IDF.
Отпечаток зависит от URL, заголовков и текстов; он позволяет
проверять согласованность индексов одного корпуса. Это служебный
некриптографический отпечаток, не доказательство отсутствия коллизий.

\journal
Пары булевой индексации расширены координатами. Добавлены внешняя
сортировка позиций, агрегирование по терму и документу, курсор
постинг-листа и проверка цитат. Разбор кавычек и \code{/ N}
встроен в общий парсер. Дистанция сравнивается в расширенном
целом типе, чтобы сумма позиции и длины фразы не переполнялась.

\testing
Для документов \code{alpha beta beta gamma}, \code{alpha x beta gamma},
\code{beta alpha gamma} проверяются точный порядок и разрешённые
вкрапления. Запрос \code{"alpha beta"} не должен находить второй
документ, а \code{"alpha beta" / 2} должен. Повторная цитата
\code{"alpha alpha"} проверяет разные вхождения. Незакрытая
кавычка и отрицательная дистанция должны вызывать ошибку.
Все четыре теста координатного набора повторно пройдены.

\results
Для полного корпуса получены 20\,430\,767 термов,
640\,214\,849 пар «терм, документ» и 2\,483\,697\,508 позиций.
Среднее число позиций на терм - 121,5665, на постинг - 3,8795.
Вариант raw без прыжков занимает 16\,694\,523\,464 байта,
координатная часть - 15\,056\,508\,824 байта. Его
сборка заняла 2214,25 секунды (36,90 минуты), создала 29 блоков,
скорость составила 7601,84 KiB/с. Пиковый RSS составил
1\,230\,024 KiB, около 1201,20 MiB. Эти времена относятся именно
к raw-сборке из эксперимента №11.

\begin{table}[H]\centering
\begin{tabular}{lrr}\toprule
Запрос & Совпадений & Время, мс\\\midrule
\code{memory allocation} & 36\,917 & 233,170\\
\code{mutex locking} & 3\,230 & 24,931\\
\code{workqueues} & 43 & 0,106\\
\code{mprotect} & 621 & 0,325\\
\code{file descriptors} & 16\,623 & 234,695\\
\code{kernel sanitizers} & 35 & 3,519\\
\code{epolll} & 0 & 0,067\\
\bottomrule\end{tabular}
\caption{Одиночные поисковые замеры координатного движка по полному корпусу; это не среднее нескольких повторов.}
\end{table}

В таблице приведены обычные строгие запросы координатного движка.
Запрос \code{"memory allocation"} на полном корпусе находит
16\,449 документов, а \code{"memory allocation" / 5} находит
17\,815 документов. Расширение дистанции допускает вкрапления
между словами и увеличивает множество совпадений.
Проверены совпадение выдачи raw- и VB-индексов, вариантов с разным
шагом прыжков, а также поиск без доступа к исходному корпусу.
Цитата может сочетаться с фильтром, например
\code{"memory allocation" && !"out of memory"}.

\conclusions
Координаты позволяют отличить совместное присутствие слов от фразы.
Цена - значительно больший объём: сохраняются все вхождения.
Частые слова и длинные списки позиций создают дорогие запросы.
Для ускорения полезны ведущий редкий терм, отбрасывание неподходящих
документов до проверки координат, сжатие и прыжки. «Точная цитата»
здесь точна по нормализованным токенам, а не по исходной строке:
регистр, пунктуация и словоформы могут различаться.

\lab{11}{Сжатие индекса}
\task
Применить сжатие к координатным блокам, описать битовое представление,
исследовать размеры и скорость прохода и поиска. Проверить сохранение
результатов и качества при преобразовании без потерь.

\method
Выбран unsigned LEB128: младшие 7 бит байта содержат очередную группу
разрядов числа, старший бит равен 1 при наличии следующего байта.
Группы записываются от младших разрядов к старшим. Для 32-битового числа
достаточно от одного до пяти байтов.
\begin{table}[H]\centering
\begin{tabular}{rll}\toprule
Число & Байты, hex & Биты\\\midrule
0 & \code{00} & \code{00000000}\\
127 & \code{7F} & \code{01111111}\\
128 & \code{80 01} & \code{10000000 00000001}\\
16383 & \code{FF 7F} & \code{11111111 01111111}\\\bottomrule
\end{tabular}
\caption{Кодирование целых LEB128.}
\end{table}
Вместо абсолютных \(docID\) записываются разности с предыдущим
документом. Позиции также кодируются разностями, но предыдущая позиция
сбрасывается в начале каждого документа. Частота вхождений остаётся
обычным беззнаковым числом, закодированным LEB128. Нулевой первый
идентификатор и нулевая первая позиция допустимы.

В raw-варианте хранятся абсолютные \(docID\) и позиции по 4 байта.
Таким образом, сравниваются практические представления raw и
«разности + VB», а не только замена кодека при одинаковых числах.
Для \(P\) постингов и \(T\) позиций raw-координаты занимают
\(4(2P+T)\) байт, что позволяет независимо сверить размер секции.

\begin{lstlisting}[language=C++]
do {
    unsigned char byte = value & 127;
    value >>= 7;
    if (value) byte |= 128;
    require(fwrite(&byte, 1, 1, file) == 1, "Write error");
} while (value);
\end{lstlisting}
Декодер контролирует конец секции, переполнение пятого байта
и неканоническое представление с лишним нулевым старшим блоком.

\journal
Чтение и запись целых вынесены в общий кодек, формат помечен флагом
сжатия. При построении сохраняются те же термы, частоты и метаданные.
В эксперименте прыжки отключены, чтобы изменение их таблицы не
влияло на сравнение размеров. Поисковые запросы выполняются после
одного прогревочного прохода, затем повторяются семь раз.
Построение и измерение разделены: измеритель читает готовые raw-
и VB-индексы и не требует исходного корпуса. Повторная подготовка
raw-индекса при существующем файле отклоняется без удаления
индекса и других результатов эксперимента.

\testing
Кодирование и декодирование проверяются на 0, границах 7-битовых
групп, \code{UINT32_MAX} и 1000 воспроизводимых псевдослучайных числах.
Проверяются некорректные байты и обрыв данных. Поиск по raw, VB
и индексу с многоблочной сборкой должен возвращать один эталонный
список. Четыре теста повторно пройдены. В полном эксперименте
дополнительно сравниваются статистика и контрольная сумма результатов.

\results
\begin{table}[H]\centering
\begin{tabular}{lrr}\toprule
Показатель & raw & VB\\\midrule
Полный индекс, байт & 16\,694\,523\,464 & 6\,510\,260\,552\\
Координатная часть, байт & 15\,056\,508\,824 & 4\,872\,245\,912\\
Прыжки, байт & 0 & 0\\
Среднее время запроса, мс & 42,546 & 48,289\\
Медианное время, мс & 17,787 & 20,540\\
Прочитано постингов & 39\,689\,783 & 39\,689\,783\\\bottomrule
\end{tabular}
\caption{Одинаковые 30 запросов, семь повторов после прогрева. Число прочитанных постингов суммарное по измеряемым повторам.}
\end{table}
Полный файл сократился до 38,996\% исходного размера (экономия
61,004\%, сжатие в 2,56 раза). Координатная часть сократилась примерно
на 67,64\%, в 3,09 раза. Среднее время выросло на 13,50\%.
Контрольные суммы результатов поиска в обоих случаях совпадают.
Числа документов, термов, постингов и позиций совпадают.

\begin{figure}[H]\centering
\begin{tikzpicture}
\begin{axis}[ybar,width=0.8\textwidth,height=6cm,
ylabel={Размер, GiB},symbolic x coords={raw,VB},xtick=data,
ymin=0,bar width=24pt,legend style={at={(0.5,1.02)},anchor=south,legend columns=2}]
\addplot coordinates {(raw,15.548),(VB,6.063)};
\addplot coordinates {(raw,14.022),(VB,4.538)};
\legend{Полный индекс,Координаты}
\end{axis}
\end{tikzpicture}
\caption{Изменение размера. Координаты входят в полный индекс; столбцы не суммируются.}
\end{figure}
\begin{table}[H]\centering\small
\begin{tabular}{lrrrrr}\toprule
Группа & Термов & raw, MiB & VB, MiB & raw, с & VB, с\\\midrule
Все & 20\,430\,767 & 14359,01 & 4646,54 & 165,469 & 176,139\\
Редкие & 17\,693\,418 & 340,32 & 145,71 & 18,408 & 18,034\\
Средние & 2\,302\,690 & 248,61 & 102,24 & 9,033 & 9,217\\
Частые & 434\,659 & 13770,08 & 4398,59 & 148,605 & 156,828\\
\bottomrule\end{tabular}
\caption{Объём координат и время последовательного прохода по группам документной частоты.}
\end{table}

Редкими считаются термы с документной частотой 1--2,
средними с частотой 3--19, частыми с частотой от 20.
Для полного прохода время выросло с 165,469 до 176,139 секунды,
на 6,45\%. У редких термов проход немного ускорился, на 2,03\%,
а у средних и частых замедлился на 2,04\% и 5,53\% соответственно.
Частые термы сжимаются плотнее: координатная часть составляет
31,94\% raw-размера против 42,82\% для редких термов.
Контрольные суммы, число постингов и число позиций совпадают
во всех группах. Время VB-сборки отдельно не измерено;
поисковое время не заменяет время построения.

\conclusions
Сжатие уменьшило объём существенно, но на данной нагрузке стоимость
декодирования оказалась выше выигрыша от уменьшения ввода-вывода.
Ускорение нельзя выводить только из размера файла. Совпадение
контрольных сумм и тестовых ответов подтверждает сохранение выдачи
на проверенных запросах. При одинаковых оценках релевантности
неизменная выдача сохраняет и значения поисковых метрик.
Возможны блочные кодеки и буферизованное
декодирование, но их выигрыш следует измерять отдельно.

\lab{12}{Прыжки по индексу}
\task
Добавить ускоренное продвижение по постинг-листам. Подобрать шаг
на обучающих запросах, проверить на других запросах, исследовать
скорость и накладной размер, убедиться в неизменности выдачи.

\method
Постинг-лист делится на блоки по \(S\) документов. Для каждого блока
хранится запись размером 24 байта: последний \(docID\) (4 байта),
предыдущий \(docID\) перед началом блока (4), смещение блока (8),
порядковый номер первого постинга (4), число постингов блока (4).
Таблица позволяет искать первый блок, последний документ которого
не меньше целевого \(docID\), бинарным поиском.

После перехода курсор читает нужный блок последовательно. Предыдущий
\(docID\) необходим для восстановления абсолютных значений из разностей:
без него произвольный вход в середину сжатого потока был бы невозможен.
Проверка порядкового номера и границ защищает от неправильного смещения.
Шаг 0 отключает таблицу и служит базовой линией.

Малый шаг сокращает последовательное чтение, но увеличивает таблицу
и число поисков по ней. Для \(L\) постингов доступ к блоку требует
порядка \(\log(L/S)\) обращений и до \(S\) последовательных чтений.
Эта оценка не включает стоимость позиционных списков и дисковых
переходов. Для коротких списков размер служебной записи может
оказаться сопоставимым с самими данными.

\journal
В индекс добавлены таблицы прыжков и метод \code{Cursor::advance}.
Для эксперимента проверяются шаги 0, 4, 16, 64, 128 и 256.
Подготовлены 15 обучающих и 15 тестовых запросов. Каждый набор
проходит один прогрев и 15 измеряемых повторов. Отдельный инструмент
меняет таблицу шага для существующих постингов, что позволяет
сравнивать один корпус и одну нормализацию. Он читает готовый VB-индекс
без прыжков и создаёт вариант с заданным положительным шагом,
сохраняя поток постингов. Результат преобразования побайтово
сравнивается с обычным построением на тестовом корпусе.

Шаг выбирается минимумом среднего времени на train. Тестовые данные
используются после фиксации выбора. Поэтому вариант, у которого
случайно лучше медиана на test, не назначается победителем задним числом.
Подготовка шести вариантов выполняется отдельной командой
\code{just prepare-tune}. Команда \code{just tune} только измеряет
готовые индексы. До записи отчётов она проверяет наличие всех
вариантов, кодек, нормализацию и шаг. Исходный корпус для замеров
не нужен, существующие индексы не перестраиваются.

\testing
Тесты проверяют переход к существующему и отсутствующему документу,
за конец списка, к границе блока и в середину VB-потока. Одни и те же
булевы запросы и цитаты выполняются с прыжками и без них на raw,
VB и многоблочном индексе. Четыре модульных теста повторно пройдены.
В полном опыте контрольные суммы должны совпадать отдельно внутри
train и test для всех шагов.

\results
\begin{table}[H]\centering
\begin{tabular}{rrrrr}\toprule
Шаг & train, мс & test, мс & Постингов test & Прыжки, MiB\\\midrule
0 & 67,316 & 31,027 & 28\,832\,670 & 0,00\\
4 & 181,180 & 75,950 & 12\,635\,580 & 3975,52\\
16 & 166,591 & 71,556 & 18\,184\,740 & 1332,07\\
64 & 153,841 & 66,239 & 23\,225\,865 & 679,69\\
128 & 147,579 & 63,039 & 24\,700\,275 & 572,43\\
256 & 139,670 & 59,450 & 25\,719\,630 & 519,30\\
\bottomrule\end{tabular}
\caption{Время, суммарное чтение за 15 повторов и размер таблиц прыжков.}
\end{table}

\begin{figure}[H]\centering
\begin{tikzpicture}
\begin{axis}[width=0.95\textwidth,height=6.8cm,
xlabel={Шаг прыжков; 0 - без прыжков},ylabel={Среднее время, мс},
xtick={0,1,2,3,4,5},xticklabels={0,4,16,64,128,256},grid=major,
legend style={at={(0.98,0.98)},anchor=north east}]
\addplot[blue,mark=*] table[x=i,y=train]{figures/skips.dat};
\addlegendentry{train}
\addplot[red,mark=square*] table[x=i,y=test]{figures/skips.dat};
\addlegendentry{test}
\end{axis}
\end{tikzpicture}
\caption{Среднее время поиска при 15 повторах после прогрева.}
\end{figure}
На train лучший шаг - 0. Следовательно, выбранный вариант на test
совпадает с базовым, а итоговое ускорение равно 1,00.
При шаге 4 число прочитанных постингов на test сокращается
с 28\,832\,670 до 12\,635\,580, то есть примерно на 56,18\%,
но среднее время растёт с 33,326 до 82,122 мс. Таблица при этом
добавляет 4\,168\,634\,952 байта.

Контрольные суммы одинаковы для всех шагов отдельно на train и test.
Счётчики чтения в таблице суммируются по 15 измеряемым повторам.
Медиана на train при шаге 64 составляет 27,746 мс против 63,181 мс
без прыжков, однако среднее ухудшается с 71,839 до 175,074 мс.
Это показывает неоднородность запросов: улучшение типичного запроса
не гарантирует улучшения среднего по всему набору.

\conclusions
Механизм прыжков корректно сокращает чтение постингов, но в этой
реализации не дал выигрыша по выбранному критерию. Вероятные причины
- поиск по таблице при частых вызовах advance, случайные обращения
и большой служебный объём. Они требуют профилирования, а не следуют
однозначно из одного замера. Возможные улучшения: хранение таблиц
в памяти, кэширование подходящего блока между вызовами,
адаптивный шаг и отключение прыжков для коротких списков.

\lab{13}{Ранжирование TF-IDF}
\task
Реализовать TF-IDF, разрешить частичное совпадение слов обычного
запроса, сохранить строгость явных булевых операций. Привести удачные
и неудачные примеры и сравнить поведение со строгим поиском.

\method
Вектор документа строится по нормализованным термам. Для ненулевой
частоты \(f_{t,d}\) используются веса:
\[
tf_{t,d}=1+\ln f_{t,d},\qquad
idf_t=1+\ln\frac{N+1}{df_t+1},\qquad
w_{t,d}=tf_{t,d}\,idf_t.
\]
Вес уникального терма запроса равен \(idf_t\). Повтор одного слова
в запросе не увеличивает его вес. Итоговая оценка - косинус:
\[
score(d,q)=\frac{\sum_{t\in q\cap d}(1+\ln f_{t,d})idf_t^2}
{\sqrt{\sum_{t\in d}((1+\ln f_{t,d})idf_t)^2}
\sqrt{\sum_{t\in q}idf_t^2}}.
\]
Норма документа включает весь его словарь и вычисляется заранее
при построении индекса. Если произведение норм равно нулю,
оценка принимается равной нулю. Термы отрицательных условий
используются как фильтр, а не как положительный сигнал сходства.

Кандидаты обычного запроса соединяются по ИЛИ, поэтому достаточно
частичного совпадения. При наличии явных операторов, скобок или
цитат запрос выполняется строго, а ранжирование меняет порядок
допустимых кандидатов. Нечёткий режим здесь означает ослабление
булевого условия, а не исправление опечаток. Неизвестные термы
не создают постингов и не включаются в используемый кодом набор
найденных термов запроса.

\begin{lstlisting}[language=C++]
double idf = log((double(index.header.documents) + 1)
                 / (term.documents + 1)) + 1;
query_norm += idf * idf;
// Contribution for a matched posting:
hits[j].score += (1 + log(double(cursor.positions.size)))
                 * idf * idf;
\end{lstlisting}
После нормировки результаты сортируются по убыванию score,
при равенстве - по возрастанию \(docID\). Текущая реализация
собирает всех кандидатов и полностью сортирует их, даже если
выводится только пять результатов.

\journal
Координатный индекс дополнен длинами документов и нормами.
Затем добавлен сбор уникальных положительных термов запроса,
накопление скалярного произведения, нормировка и сортировка.
Строгий и ранжированный движки используют один индекс;
в конфигурации переключается \code{search_engine=search|rank}.
Для честного сравнения сохранены обе выдачи одних 30 запросов.
Измеритель открывает основной готовый индекс и проверяет соответствие
кодека, нормализации и шага конфигурации. Он не строит дополнительный
индекс, не удаляет каталог результатов и работает без исходного корпуса.

\testing
На малом корпусе норма рассчитывается вручную и сравнивается
с допуском \(10^{-12}\). Проверяются нулевая норма пустого
документа, одинаковые веса, детерминированное разрешение равенства,
строгое И и частичное совпадение. Режим rank не должен превращать
явный \code{&&} в ИЛИ. Четыре модульных теста пройдены.
CLI-проверка дополнительно выполняет измерение дважды при недоступном
корпусе, сверяет неизменность индекса и проверяет ошибку при его отсутствии.

\results
\begin{table}[H]\centering
\begin{tabular}{lrr}\toprule
Запрос & Строгий поиск & rank (ИЛИ)\\\midrule
\code{memory allocation} & 36\,917 & 282\,185\\
\code{mutex locking} & 3\,230 & 30\,501\\
\code{waiting children} & 7\,543 & 216\,266\\
\code{file descriptors} & 16\,623 & 753\,972\\
\code{memory copies} & 85\,345 & 545\,106\\
\code{epolll} & 0 & 0\\
\bottomrule\end{tabular}
\caption{Число кандидатов в выдачах строгого и ранжированного поиска.}
\end{table}

Для запроса \code{memory allocation} первым идёт \emph{Page Allocation}
с оценкой 0,749056, затем \emph{Virtually Contiguous Memory Allocation}
с 0,539702 и \emph{C++ keyword: new} с 0,455739.
Однако первые две страницы в корпусе содержат только краткие
заголовки, а третья описывает ключевое слово C++. Поэтому все три
получили оценку релевантности 1 для запроса об инструкциях выделения
памяти в Linux и libc. Тематически подходящий заголовок ещё не означает
наличия ответа в документе. Строгий поиск начинает с
\emph{AddressSanitizer}, который также не отвечает непосредственно на запрос.

Удачный пример даёт запрос \code{malloc}: первыми выводятся
\emph{Unconstrained Allocation} и \emph{Basic Allocation} из
документации GNU libc. Для этого запроса NDCG@5 равна 1,0000.

Есть и менее удачные примеры. Для \code{memory copies} и
\code{memory overlapping} первым становится \emph{Out Of Memory
Handling}: частичное совпадение по memory не отражает намерение
найти копирование перекрывающихся областей. Для \code{printf}
первой идёт задача \emph{Number triangle with while loop in c},
а не справочная страница функции. На \code{epolll} результатов нет.
В этих случаях нужны соответственно учёт связи слов, предпочтение
API-страниц и отдельная коррекция опечаток.

Бенчмарк 30 запросов после одного прогрева и с семью повторами:
среднее 806,199 мс, медиана 264,808 мс, суммарный счётчик прочитанных
постингов 79\,385\,544. Это существенно дороже строгого
поиска, поскольку обычные запросы создают широкое объединение
кандидатов и требуют вычисления весов. Замеры разных лабораторных
не дают чистой оценки цены одной операции: у них отличаются
наборы параметров и путь исполнения.

По полностью размеченному пулу из лабораторной №6 получены
P@5 = 0,4933, DCG@5 = 7,8437, NDCG@5 = 0,4219,
ERR@5 = 0,4742. Для строгого поиска значения равны соответственно
0,1933, 2,9820, 0,1744 и 0,2283.
Сравнение оценивает весь переход от строгого поиска с порядком
по идентификатору к объединению кандидатов и TF-IDF, а не только
влияние одной формулы весов. Рост числа кандидатов сам по себе
не является улучшением качества.

\conclusions
TF-IDF выводит сначала документы с более сильным лексическим
совпадением, сохраняя строгие фильтры. Основные недостатки -
независимость термов, отсутствие смыслового понимания и высокая
цена полного ранжирования широкого множества. Улучшения:
top-k-отбор вместо полной сортировки, BM25, отдельный вес заголовка,
бонус близости слов и защита технических имён. Каждое изменение
следует проверять на одном размеченном пуле.

\lab{17}{Автотесты (юнит-тесты)}
\task
Подготовить воспроизводимые модульные тесты основных компонент
поисковой системы. Проверить штатные и граничные случаи,
некорректные входы, согласованность форматов и результатов.

\method
Используется Google Test, вызывающий функции C++ напрямую.
CMake подключает установленный пакет или закреплённую версию
1.15.2 с проверкой SHA-256. CTest регистрирует отдельные тесты,
поэтому отчёт показывает конкретную проверку, вызвавшую ошибку.
Тесты запускаются отдельно от обычной сборки.

Искусственный корпус содержит восемь документов:
\begin{center}
\begin{tabular}{rl}\toprule
ID & Текст\\\midrule
0 & \code{alpha beta beta gamma}\\
1 & \code{alpha x beta gamma}\\
2 & \code{beta alpha gamma}\\
3 & \code{alpha beta gamma}\\
4 & пустой документ\\
5 & \code{delta gamma}\\
6 & \code{alpha alpha beta}\\
7 & \code{common common delta}\\\bottomrule
\end{tabular}
\end{center}
Ручной эталон фиксирует ответ независимо от поискового кода.
Например, \code{"alpha beta"} даёт \(\{0,3,6\}\),
\code{"alpha beta" / 2} - \(\{0,1,3,6\}\),
\code{!alpha} - \(\{4,5,7\}\).
Один набор запросов выполняется по raw, VB и индексу,
построенному с очень малым буфером, с прыжками и без них.

\journal
Сначала добавлены тесты кодека и булевых выражений. Затем они
расширены цитатами, внешним слиянием, метаданными и нормами TF-IDF.
Для каждого случая используется временный каталог, а данные
создаются самим тестом: рабочий корпус и готовые индексы
не перезаписываются. \code{SCOPED_TRACE} добавляет к диагностике
текущий запрос. Для псевдослучайных целых кодека задано фиксированное зерно.

\testing
\begin{longtable}{p{0.41\textwidth}p{0.51\textwidth}}
\toprule Проверка & Что обнаруживает\\\midrule
\code{BoundariesAndInvalidEncoding} & Ошибки границ LEB128, каноничности и восстановления целых.\\
\code{BooleanAndOrderedPhrases} & Ошибки приоритетов, цитат, порядка, отрицания и повторных слов.\\
\code{MergeMetadataNormsAndSkips} & Зависимость файла от размера блока, повреждение метаданных, норм и курсоров.\\
\code{FuzzySearchAndTfIdfRanking} & Нарушение строгих фильтров, частичных совпадений и ранжирования.\\\bottomrule
\end{longtable}
Ручная норма первого документа использует
\(idf=1+\ln(9/6)\), \(\beta=1+\ln2\),
\(\|d_0\|=idf\sqrt{2+\beta^2}\).
Пустой документ имеет нулевые длину и норму. Побайтовое
сравнение VB-файлов при разном размере буфера проверяет,
что внешнее слияние не меняет результат сборки.

\results
Повторно запущены модульные тесты и проверки командного интерфейса.
Тесты создают только небольшие временные данные, полные индексы
не перестраиваются. Получены результаты:
\begin{table}[H]\centering
\begin{tabular}{lrr}\toprule
Проект & Пройдено & Ошибок\\\midrule
Оценка качества, №6 & 3 & 0\\
Булев индекс, №7 & 2 & 0\\
Булев поиск, №8 & 2 & 0\\
Координатный поиск, №10 & 4 & 0\\
Сжатие, №11 & 4 & 0\\
Прыжки, №12 & 4 & 0\\
TF-IDF, №13 & 4 & 0\\
Общий набор, №17 & 4 & 0\\
Проверка оценочного конвейера & 7 & 0\\
Регрессионные проверки CLI & 6 & 0\\\midrule
Итого запущено & 40 & 0\\\bottomrule
\end{tabular}
\caption{Результаты C++- и Python-тестов. Наборы 10--13 и 17 частично дублируются.}
\end{table}
Фрагмент журнала общего набора:
\begin{lstlisting}
[ RUN      ] Codec.BoundariesAndInvalidEncoding
[       OK ] Codec.BoundariesAndInvalidEncoding
[ RUN      ] PositionalIndex.BooleanAndOrderedPhrases
[       OK ] PositionalIndex.BooleanAndOrderedPhrases
[ RUN      ] PositionalIndex.MergeMetadataNormsAndSkips
[       OK ] PositionalIndex.MergeMetadataNormsAndSkips
[ RUN      ] PositionalIndex.FuzzySearchAndTfIdfRanking
[       OK ] PositionalIndex.FuzzySearchAndTfIdfRanking
[  PASSED  ] 4 tests.
\end{lstlisting}
Дополнительные семь тестов оценивания проверяют валидность
и полноту разметки. Шесть регрессионных CLI-тестов проверяют
повторные измерения без корпуса, неизменность индексов,
ошибки при отсутствии подготовленных файлов, отказ от повторной
записи существующего индекса и точность преобразования шага прыжков.
Для сравнения индексов используются SHA-256 и побайтовая проверка.

Число 40 означает количество выполненных тестов, а не число
независимых алгоритмов: часть C++-наборов повторяется в разных проектах.
Автотесты робота, Python-обвязки токенизатора и нормализатора
в этот итог не включены.

\conclusions
Малый корпус с ручными ответами позволяет проверять семантику
и ловить регрессии при смене кодека и способа сборки. Тесты
не доказывают отсутствие ошибок и не являются бенчмарком.
HTTP-интерфейс и настоящая MongoDB требуют отдельных интеграционных
проверок. Длительные сбои диска и сети требуют нагрузочных и аварийных
испытаний. Релевантность оценивается на размеченном пуле, а не
модульными тестами алгоритма.

\clearpage
\section*{Сборка и воспроизведение}
\addcontentsline{toc}{section}{Сборка и воспроизведение}
Для C++ нужны компилятор с C++20, CMake 3.24+ и \code{just}.
Для Python используются \code{uv}, файлы \code{pyproject.toml}
и \code{uv.lock} каждого проекта. Для живого робота нужны Docker
и MongoDB replica set. Первая настройка тестов может требовать
сети для загрузки закреплённого Google Test.

Обычная сборка \code{just build} не строит индекс.
Команда \code{just build-index} - отдельная дорогостоящая
операция; для уже сохранённого индекса она не нужна.
Относительные пути в конфигурации считаются от каталога проекта.
Пример поиска:
\begin{lstlisting}
cd 13_tfidf
just build
printf 'memory allocation\n"memory allocation" / 5\n' |
  just search
just web
\end{lstlisting}
Команда тестирования:
\begin{lstlisting}
cd 17_unit_tests
just test
ctest --test-dir build --output-on-failure
\end{lstlisting}
Первая команда собирает и запускает тесты, вторая повторяет их
без пересборки. Для всех C++-проектов принцип тот же.

Измерение сжатия в лабораторной №11 выполняется после отдельного
построения VB- и raw-индексов командами \code{just build-index}
и \code{just build-raw}, затем \code{just measure}.
При уже существующем raw-индексе повторная подготовка завершается
ошибкой и не удаляет готовые файлы.

В лабораторной №12 команда \code{just prepare-tune} подготавливает
шесть вариантов индекса. Затем \code{just tune} измеряет их и
выбирает шаг. Для измерений корпус не нужен; отсутствие варианта
вызывает ошибку, а не автоматическую перестройку.
В лабораторной №13 \code{just measure} использует основной готовый
индекс и сравнивает строгий и ранжированный поиск.

Для пересчёта качества из каталога лабораторной №6:
\begin{lstlisting}
just evaluate-full
just plot-full
\end{lstlisting}
Первая команда проверяет согласованность запросов, выдач и оценок,
а затем вычисляет метрики. Вторая дополнительно строит графики.
Для этих команд не требуются корпус и поисковые индексы:
они работают по зафиксированным выдачам и разметке.
Неоценённый результат или изменение проверенных входных данных
приводит к диагностической ошибке.

\clearpage
\subsection*{Запросы и информационные потребности}
\begin{longtable}{lp{0.28\textwidth}p{0.51\textwidth}}
\toprule ID & Запрос & Информационная потребность\\\midrule\endfirsthead
\toprule ID & Запрос & Информационная потребность\\\midrule\endhead
q01 & \code{memory allocation} & Как выделять память в Linux и libc?\\
q02 & \code{memory allocations} & Тот же запрос во множественном числе; устойчивость к словоформе.\\
q03 & \code{mutex locking} & Как блокировать mutex в ядре и pthreads?\\
q04 & \code{spinlocks} & Как работают spinlock и когда их применять?\\
q05 & \code{workqueues} & Как планировать отложенную работу через workqueue?\\
q06 & \code{reference counting} & Как вести счётчик ссылок и не допустить use-after-free?\\
q07 & \code{memory mappings} & Как отображать и снимать отображение памяти?\\
q08 & \code{mprotect} & Как изменить права доступа к страницам памяти?\\
q09 & \code{fork} & Как создать дочерний процесс через fork?\\
q10 & \code{execve} & Как заменить образ процесса через execve?\\
q11 & \code{waiting children} & Как дождаться завершения дочерних процессов?\\
q12 & \code{pipes} & Как передавать данные через pipe?\\
q13 & \code{file descriptors} & Как открывать, дублировать и обслуживать файловые дескрипторы?\\
q14 & \code{synchronous writes} & Как гарантировать запись данных на носитель?\\
q15 & \code{epoll events} & Как получать уведомления готовности epoll?\\
q16 & \code{sockets} & Как создать и использовать сетевые сокеты?\\
q17 & \code{accept connections} & Как принимать входящие соединения?\\
q18 & \code{futex} & Как устроена синхронизация через futex?\\
q19 & \code{CPU affinity} & Как привязать задачу или прерывание к CPU?\\
q20 & \code{resource limits} & Как установить ограничения ресурсов процесса?\\
q21 & \code{pthread_create} & Как создать POSIX-поток? Точное имя API.\\
q22 & \code{malloc} & Как выделять память через malloc? Точное имя API.\\
q23 & \code{memory copies} & Как копировать блоки памяти?\\
q24 & \code{memory overlapping} & Как безопасно копировать перекрывающиеся области памяти?\\
q25 & \code{printf} & Как форматировать вывод? Точное имя API.\\
q26 & \code{dlopen} & Как загрузить разделяемую библиотеку? Точное имя API.\\
q27 & \code{huge pages} & Как использовать huge pages?\\
q28 & \code{cgroup memory} & Как ограничить память через cgroup v2?\\
q29 & \code{kernel sanitizers} & Как обнаружить ошибки памяти и гонки в ядре?\\
q30 & \code{epolll} & Опечатка в epoll; ожидается пустая выдача без коррекции.\\
\bottomrule
\end{longtable}


\clearpage
\addcontentsline{toc}{section}{Список литературы}
\begin{thebibliography}{9}
\bibitem{course} Кухтичев А.\,А. Материалы курса «Информационный поиск»,
МАИ, осень 2025: задания, лекции и шаблон отчёта.
Репозиторий MAI-IR: условия лабораторных работ,
лекционные материалы и домашние задания лекций 1, 2, 4.
\bibitem{manning} Маннинг К.\,Д., Рагхаван П., Шютце Х.
Введение в информационный поиск. М.: Вильямс, 2011. 528 с.
\bibitem{dataset} Common Pile. Stack Exchange: архивный корпус.
\url{https://huggingface.co/datasets/common-pile/stackexchange}.
Использованная ревизия указана в лабораторной №1.
\bibitem{wordnet} Princeton University. WordNet 3.0.
Словарные данные и исключения, включённые в исходники проекта;
лицензия сохранена вместе со словарём.
\bibitem{porter} Porter M.\,F. An algorithm for suffix stripping.
Program. 1980. Vol. 14, no. 3. P. 130--137.
\bibitem{implementation} Исходные тексты лабораторных работ №1--8,
10--13, 17. Собственные реализации индексатора, поисковых алгоритмов,
оценки качества и автоматических тестов; результаты экспериментов
по корпусу документации и обсуждений программирования.
\end{thebibliography}

\end{document}
